\chapter{Introduction}
\label{chap:intro}

Stock markets are not only trading venues. They are mechanisms for capital allocation, corporate financing, and risk transfer. When public information is absorbed into prices, the speed and quality of that absorption affect investors, listed firms, and the broader economy. For research desks and brokerage institutions, the practical question is therefore not merely whether prices move, but whether new information can be interpreted quickly enough to support better decisions before the reaction is fully reflected in the market. This economic relevance is one reason why stock forecasting remains a persistent research problem. \cite{fama1970, lo1988randomwalk, sezer2020dlreview}

The technical difficulty begins with an information asymmetry. Numerical market series record the market reaction after it occurs, whereas news, disclosures, and commentary describe the events that may have triggered that reaction. A price chart reveals that a move happened. It does not explain whether the move followed earnings news, financing plans, governance changes, or contract announcements. This gap motivates the use of text as an exogenous signal rather than relying on price history alone. \cite{tetlock2007media, kearney2014survey, sezer2020dlreview}

Existing research has therefore moved from price-only forecasting toward multimodal forecasting that combines market variables with textual sentiment, event representations, or contextual language-model features. The evidence is encouraging but not uniform. Some studies report measurable gains from financial news or social-media sentiment, while others show that those gains depend strongly on the text source, the forecast horizon, the market under study, and the evaluation protocol. The relevant scientific question is conditional: under what data and validation conditions does sentiment add predictive value? \cite{das2007yahoo, bollen2011twitter, dahal2023comparative, liapis2023deep}

This question is especially important in the Egyptian market, where firm-related financial coverage appears in both Arabic and English. A forecasting pipeline that ignores one of the two languages discards part of the information environment. A pipeline that combines both languages without careful routing and relevance control can instead amplify noise. The present thesis emerged from an industrial bachelor collaboration with Robin, the data-science arm of Beltone Holding, where this issue was practical rather than abstract. The working need was to test whether bilingual financial-news sentiment could support EGX-oriented forecasting analysis under a setting that remained interpretable, conservative, and reproducible. Prior work on finance-specific English sentiment modeling, Arabic-specialized language models, and conservative sentiment-evaluation design supports that emphasis on language-aware modeling and strict temporal alignment. \cite{araci2019finbert, antoun2020arabert, abdulmageed2021arbert, liapis2023deep}

\section{Problem Statement}

This thesis investigates whether sentiment extracted from bilingual financial news can improve short-horizon stock direction forecasting in order to determine the predictive value of sentiment-derived signals for selected Egyptian Exchange equities, namely CIB, TMG, and OIH.

\section{Motivation}

The industrial context of the study is a strong motivation for the chosen design. In the Robin and Beltone workflow that motivated the thesis, analysts face a mixed Arabic-English news stream, uneven company coverage, and many articles that are attached to a company page without being equally relevant to valuation. A useful forecasting signal in that setting must therefore do more than assign positive or negative labels. It must route each language to an appropriate sentiment model, distinguish direct company coverage from weaker mentions, and aggregate articles in a way that preserves timing and relevance. This strategy is also consistent with prior work showing that source choice, finance-specific English sentiment modeling, and Arabic-specialized encoders materially affect the quality of textual sentiment signals. \cite{kearney2014survey, araci2019finbert, antoun2020arabert}

The second motivation is methodological. A strong price-based forecasting system can hide whether sentiment is truly informative or whether it is merely a small additive feature inside an already expressive model. For that reason, the thesis adopts a sentiment-only benchmark before any later integration with richer market features. This choice makes the interpretation cleaner and aligns with the cautionary literature on sentiment forecasting, where apparent gains often disappear once feature leakage, horizon choice, or loose validation procedures are controlled more strictly. \cite{liapis2023deep, dahal2023comparative, sezer2020dlreview}

\section{Research Questions}

\noindent\textbf{RQ1.} To what extent does bilingual financial-news sentiment contain predictive information for short-horizon EGX stock direction forecasting?

\noindent\textbf{RQ2.} How do language-specific routing, article relevance, and event-aware aggregation affect the quality of the extracted sentiment signal?

\noindent\textbf{RQ3.} Do any observed predictive gains remain valid under strictly chronological and leakage-safe evaluation?

\noindent\textbf{RQ4.} What do auxiliary experiments with historical-model integration, sentiment-return correlation, Arabic financial-domain adaptation, and social-media feasibility reveal about the future development path of EGX sentiment forecasting?

\section{Objectives}

\noindent\textbf{Objective 1.} Design a bilingual financial-news forecasting formulation for selected EGX equities.

\noindent\textbf{Objective 2.} Develop a sentiment extraction framework that distinguishes Arabic and English article processing.

\noindent\textbf{Objective 3.} Construct sentiment features that account for article relevance, language composition, and event structure.

\noindent\textbf{Objective 4.} Evaluate the resulting forecasting signal under chronological validation and test splits with explicit leakage controls.

\noindent\textbf{Objective 5.} Investigate whether the observed sentiment signal is strong enough to justify further sentiment-enhanced forecasting work in the Egyptian market.

\noindent\textbf{Objective 6.} Document auxiliary experiments that test sentiment as an added feature in a historical model, inspect direct sentiment-return association, evaluate Arabic financial-domain adaptation, and define ethical future work for social-media sources.

\section{Thesis Outline}

The remainder of this thesis is organized as follows. Chapter~\ref{chap:background} establishes the background concepts required for the study, including forecasting targets, class imbalance in directional prediction, financial text as an exogenous information source, transformer-based sentiment extraction, and leakage-safe evaluation. Chapter~\ref{chap:related} then reviews state-of-the-art work on sentiment-based stock forecasting and positions the present thesis within that literature. Chapter~\ref{chap:methodology} presents the solution approach adopted in the study. Chapter~\ref{chap:results} reports the implementation choices, experimental setup, and empirical findings. Finally, Chapter~\ref{chap:conclusion} summarizes the main conclusions and outlines future research directions.
